\subsection{Overview}

Watershed modeling occupies a pivotal but unfinished position in the science of drainage water management: it is the principal tool for translating field-scale process understanding into landscape-scale predictions of nutrient and sediment fate, yet the models themselves, the data that constrain them, and the scales at which they are applied remain the subject of active methodological debate. The recurring call across this literature is to \emph{refine, build, calibrate, and validate} watershed hydrological models used to assess how drainage water management affects the nitrogen and phosphorus cycles \parencite{Sunohara2015}, and to pair such modeling with large-scale efforts that prioritize where drainage water management will yield the greatest benefit \parencite{Ross2016}. Several intertwined research lines emerge from the findings. First, models require improved representation of processes that are currently poorly resolved---phosphorus sorption--desorption, preferential flow, in-stream processing, channel-network routing, and reactive transport---together with explicit treatment of uncertainty \parencite{Wang2020b,Kujawa2020}. Second, a persistent scale problem separates well-characterized fields from the watersheds where water-quality goals are set, motivating explicit field-to-watershed linkage and upscaling frameworks \parencite{Williams2015,Kyllmar2023}. Third, the empirical base for calibration and validation remains thin: long-term, year-round, coordinated multi-site monitoring is repeatedly identified as the binding constraint \parencite{Christianson2015,Bauwe2020}. Fourth, the field is called on to expand beyond its intensively tile-drained Midwestern core to less intensely drained, newly drained, and geographically underrepresented watersheds \parencite{Kujawa2020,Christianson2015}. Finally, modeling is increasingly asked to serve decision support---placement of stacked conservation practices, adaptive management, and prioritization of critical source areas---rather than description alone \parencite{Kröger2013,Tomer2015}. The subsections below organize these open questions by the modeling task they most directly challenge.

\subsection{Refining process representation in watershed models}

A large share of the frontier concerns processes that current watershed models represent crudely or omit. For phosphorus, \textcite{Wang2020b} set out an explicit development agenda: incorporate dynamic P sorption/desorption processes and new subroutines for direct P loss from manure, fertilizer, and dung; integrate pedotransfer functions for the P sorption coefficient and distinct P pools keyed to soil characteristics and weathering; include in-stream phosphorus processes and critical-source-area identification based on variable source areas; and, crucially, model preferential flow and connectivity between fields and adjacent water bodies. Preferential flow recurs as a first-order gap. \textcite{Djodjic2002} argue that preferential-flow losses of P in structured soils must be incorporated into phosphorus indices and decision-support systems, that identification of preferential pathways at the field scale would sharpen leaching-risk assessment, and that employing more than one hydrological model within a decision-support system may better describe water and pollutant transport. Consideration of preferential flow in the determination of critical source areas, especially at field or sub-field scale, is likewise flagged as necessary \parencite{Wang2020b}. The adequacy of our understanding of subsurface hydrologic pathways---and whether they can be accurately modeled at the watershed scale---remains an open question in its own right \parencite{Sims1998}.

Channel networks and in-stream processing are a second underdeveloped domain. \textcite{Giraud1997} emphasize that as watershed size increases, channel networks must be modeled as a whole set of interdependent reaches---whether looped or dendritic---and that modeling flow paths and travel time of water within the network is essential to predicting non-point-source pollution reaching coastal waters. At the watershed scale, more detailed methods of channel-network description and extrapolation may reduce uncertainty, and evaluation of historical land management and channel stability from aerial photography could further inform the process \parencite{Williamson2025}. Better representation of instream biogeochemical processes is proposed as a route to enhance the reductions generated by best management practices \parencite{Schilling2012b}, and the interaction between hillslopes and riparian areas needs stronger representation to predict low-flow events \parencite{Basu2010}.

Sediment dynamics require analogous refinement. Estimating gully and river-bed erosion will demand different approaches than upland modeling, and the contribution of gully areas to the overall sediment budget deserves closer scrutiny, while the use of WEPP for upland sediment flux needs verification against well-recorded observations \parencite{Maalim2013}. A more robust treatment of sediment from non-field sources should explicitly represent the mechanisms and physical processes behind streambank failure and downstream mobilization \parencite{Dalzell2018}. At the catchment scale, sediment budgets could underpin the design of control and management strategies for lowland agricultural catchments with extensive underdraining, and their precision should be carried forward explicitly into future budget refinement \parencite{Walling2002}.

\subsection{Uncertainty, model evaluation, and structural choices}

Beyond adding processes, the findings press for rigorous quantification of uncertainty and disciplined model evaluation. Assessments of model uncertainty should be built into P-transport modeling \parencite{Wang2020b}, and structural uncertainty deserves explicit treatment---for example by including water-quality models beyond SWAT when projecting hydrologic and nutrient change \parencite{Kujawa2020}. Standardized evaluation metrics recur as a methodological aspiration, with the coefficient of efficiency of Nash and Sutcliffe proposed for future watershed-model evaluations \parencite{STONE1995,SAVABI1995}. For the specific problem of drainage water management effectiveness, \textcite{Shedekar2021} argue that a probabilistic approach accounting for stochastic variation in effectiveness across time and scales may be more appropriate than a single static value or range, and that large datasets combining short- and long-term monitoring and modeling at regional scales may soon be needed to support large-scale assessments. Model transferability is itself a research question: whether a representative 10--20~km\textsuperscript{2} scale identified in one watershed applies to others in the region needs to be examined, and the validity of a given modeling approach must be established by comparison against hydrograph data from additional watersheds \parencite{Basu2010}. Newly proposed physically based tools warrant similar scrutiny---for instance, applying the Advanced Terrestrial Simulator (ATS) to larger watersheds with different soils and topography to probe scaling and physiographic controls, and extending it to explore land-use change, surface-drainage modification, future climate, and drainage intensification \parencite{Rathore2025}. Where field data are unavailable, calibrated and validated physically based models such as HYDRUS-2D are recommended for reliable flux estimation, provided calibration and validation datasets are made more comprehensive \parencite{Naderi2025}. Individualized, station-by-station analysis is also advised for regional trend work, since a single approach may not yield accurate results across heterogeneous stations \parencite{Kumar2009}.

\subsection{The field-to-watershed scaling problem}

A defining open problem is the disconnect between field-scale process knowledge and watershed-scale response. \textcite{Williams2015} identify quantifying the relationships between field-scale processes and watershed response as a critical need for developing water-quality strategies, and cast the field--watershed linkage as pivotal for tile-drained regions of the Midwest and beyond. Integrating field and watershed models for BMP calibration---and scaling field results up to larger areas---is an explicit modeling priority \parencite{Wang2020b}. Small, intensively monitored catchments are proposed as the empirical bridge: they can reveal system functioning, serve as a basis for upscaling to agricultural land where information is less detailed, and provide the data for calibrating and validating modeled outputs \parencite{Kyllmar2023}. Scaling questions also arise for individual practices. Effectiveness estimates for drainage water management should be scaled up from field to watershed \parencite{Wang2021}, and phosphorus mitigation likewise needs scale-up from field to watershed in future work \parencite{Wang2021}. Selenium-removal and other treatment findings are called to be extended to scaled-up operational flow-through systems \parencite{Gao2003}, and complex questions remain about the effectiveness of widespread practice implementation at watershed scales, compounded by lag times in water-quality response and low adoption rates in key watershed locations \parencite{Roley2016}. Continued large-scale watershed modeling is repeatedly recommended to prioritize the areas that will benefit most from drainage water management \parencite{Ross2016}.

\subsection{Data, monitoring, and long-term records to constrain models}

Nearly every modeling ambition is bottlenecked by data. \textcite{Christianson2015} call for more intensive year-round monitoring in some locations, improved monitoring of newly drained areas, ditch-drained areas, and surface intakes, more long-term studies with coordinated controls across multiple sites and years, and continued construction of region-specific drainage water-quality data for underrepresented regions---adding that understanding drainage-design impacts grows more important as aging infrastructure is replaced and upgraded. Long-term records are singled out as necessary to disentangle the possibly counteracting effects of land management and climate change on biogeochemical and hydrological processes, with the corollary that evaluating policy measures will demand patience \parencite{Bauwe2020}. Long-term, catchment-scale study is deemed essential to understand management and climate-change impacts \parencite{Kelley2017}, and the establishment and use of long-term studies tracking the effects of management changes on runoff phosphorus losses is a standing recommendation \parencite{Gillingham2000}. Follow-up investigations of implemented strategies are described as critical for adaptive management \parencite{Schilling2012b}, and long-term watershed data are advocated as the foundation for models that target high-pollution-potential areas and project watershed responses to BMP implementation \parencite{Lemke2011}.

The tracer and water-age literature adds specific data demands for constraining transport models. Further quantification of water ages in drained agricultural landscapes with high-frequency (daily to sub-daily) tracer measurements is needed, along with additional water and tracer datasets across fields with differing management, soils, and climates, and more complex models capable of representing both celerities and velocities \parencite{Williams2021}. Event sampling is flagged as an important next step to constrain the younger water fraction, and investigating the isotope dynamics of ``bulk'' soil water would more completely characterize water partitioning \parencite{DimitrovaPetrova2020}. Further quantification of water ages across drained landscapes would clarify the trade-offs between maintaining drains for food production and the exacerbation of flooding and contaminant transport \parencite{DimitrovaPetrova2020}. High-resolution specific-conductance data are expected to enable more end-member-mixing hydrograph separations as monitoring networks grow, with the caveat that method choice should reflect the flow and specific-conductance characteristics of the stream \parencite{Kronholm2015}. Improved national drainage-extent datasets open new analyses of drainage impacts and nutrient-control strategies, though further study of how farmers report ambiguous cases such as nonpatterned drainage would aid interpretation \parencite{Jame2022}. Coordinated, multidisciplinary, and international collaboration---including consistent reporting of basic performance parameters such as influent/effluent concentrations, hydraulic residence and loading rates, and watershed-to-treatment ratios---is urged to build the shared evidence base for constructed-wetland and treatment performance \parencite{Messer2021,Sims1998}.

\subsection{Extending to underrepresented and less intensely drained landscapes}

A clear frontier is geographic and physiographic breadth. Applications developed in intensely drained watersheds should be extended to watersheds that are \emph{not} intensely drained \parencite{Kujawa2020}. Regional trend and discharge analyses call for larger basin areas, longer study periods, and replication across multiple watersheds while controlling for weather variability, including the influence of small sub-field wetlands \parencite{Yang2017}. Drainage-flow analyses should be extended to larger drainage networks---for example the many US Midwest networks exceeding 29~ha---and linked to measured nitrogen loading across drainage scales and across additional states \parencite{Christianson2025}. Broader applicability is sought by explicitly addressing regional variability, external factors, and climate change \parencite{Naderi2025}, and additional studies across a broad range of climates and topographic features are recommended to strengthen emerging hydrologic paradigms for headwater streams \parencite{King2009}. The importance of physical-hydrologic dynamics as drivers of nutrient loss in dryland regions worldwide remains to be assessed \parencite{Kelley2017}. Regional hydro-climatic feedbacks---including the effects of biofuel expansion and of wetland drainage on regional hydrology and climate---are proposed as further avenues \parencite{Kumar2009}. Effects of drainage-related structures on wetland hydrology across a range of site and meteorological conditions require dedicated field data \parencite{Tangen2018}, and inland-basin remediation treatments are recommended for replication across affected areas \parencite{Kumar2006}. Theoretical drainage results are also candidates for generalization---for example to the case of deep impervious barriers and to water tables drained by channels or rivers \parencite{Zimmer2023}.

\subsection{Modeling to guide practice placement, stacking, and decision support}

Increasingly, watershed modeling is asked to inform \emph{where} and \emph{how} to deploy practices. New conservation-planning tools are needed to determine the best placement of practices in the landscape and to evaluate alternative conservation scenarios, ideally via a hydro-biogeochemically based tool that compares practice effectiveness by landscape position and identifies downstream sediment-deposition zones \parencite{Kröger2013}. Coupling practice implementation with site- and watershed-scale monitoring and modeling is presented as a way to advance water-quality improvement and practice-effectiveness assessment simultaneously, within an adaptive-management framework that maintains effectiveness over time \parencite{Kröger2013}. Simulation models are proposed to test and refine watershed scenarios, alongside continued development of placement criteria for new and emerging practices and better understanding of controls on nutrient-removal efficiency \parencite{Tomer2015}. Identifying and targeting the critical areas within watersheds that generate the largest share of nutrient and sediment export is a shared priority \parencite{Lemke2011}, echoed by calls to identify fields, farms, or watersheds exporting environmentally significant phosphorus via subsurface runoff and to improve soil-testing and characterization methods for prioritization \parencite{Sims1998}. Coarse-scale loss estimates are valued for flagging where deeper modeling and mitigation of nitrogen, phosphorus, and sediment losses should be pursued at catchment and regional scales \parencite{McDowell2021}. Decision-support systems are expected to advance through GIS-based non-point-source models that better capture spatial and temporal variability, incorporation of remote sensing, and software interfaces linking their components \parencite{Djodjic2002}. Combining ("stacking") drainage water management with other in-field and edge-of-field practices to meet phosphorus-reduction goals is an explicit design frontier \parencite{King2022}, and the nitrogen-removal capacity of stacked conservation practices is highlighted as a hypothesis to test \parencite{Roley2016}. Conservation planning in tile-drained watersheds is expected to require combining surface-water BMPs with practices that intercept and retain subsurface runoff, supported by watershed plans and by policies and incentive programs that specifically target tile-drainage nutrient sources \parencite{Lemke2011}. Wider involvement of land managers in regulatory assessments---supported by online teaching tools and analytical services---is likewise encouraged \parencite{Stutter2018}, and cost-share programs promoting drainage water management for watershed nitrogen and phosphorus targets are advocated \parencite{Sunohara2015}.

\subsection{Coupling nutrient, carbon, and greenhouse-gas processes}

The modeling frontier extends to biogeochemical coupling that current watershed tools rarely close. Reactive and sediment-transport models coupled to physically based hydrology are proposed to reveal drivers of adverse watershed impacts and to generate new insight into watershed-scale nutrient transport \parencite{Rathore2025}. For phosphorus specifically, investigating the multitude of P species and flow pathways for all relevant practices, and continuing to study surface-runoff dynamics---particularly under controlled drainage---are identified priorities \parencite{Mitchell2023}. Future work on the interaction between legacy soil phosphorus and hydrology in drained landscapes and closed depressions is needed to quantify P-loss risk \parencite{Mumbi2024}, and phosphorus-management strategies should jointly incorporate precipitation-pattern dynamics and soil-nutrient management, prioritizing application rate, timing, and placement plus legacy-P mitigation, especially where extreme weather is intensifying \parencite{Mumbi2025b}. Mechanistic controls on elevated, solubility-driven phosphorus loads to Lake Erie and other downstream waters require further study of spatial and temporal interactions between site characteristics and P solubility \parencite{Hanrahan2020}. Greenhouse-gas coupling is a notable gap: quantifying fluxes---including carbon sequestration, pre-implementation emissions, and dissolved nitrous oxide in tile drainage and conservation-practice outflows---is called for \parencite{Mitchell2023}. Coordinated, large-scale studies are needed to constrain the climate-mitigation potential of conservation drainage practices, few of which have simultaneously measured N\textsubscript{2}O, CH\textsubscript{4}, and change in soil organic carbon at the same sites, or included dissolved-gas measurements \parencite{Hall2025}. Practice-specific unknowns include the fate of water held back by control structures, interactions among bioreactor design, catchment hydrology, and organic media, controls on soil-organic-carbon accrual and N\textsubscript{2}O yield in saturated buffers, and N\textsubscript{2}O implications of drainage water recycling \parencite{Hall2025}. In the Midwest, longer-term dynamics of water-quality-wetland carbon storage, nitrous oxide, and methane over a wetland's expected lifetime---and the effects of upgrading drainage infrastructure and drainage coefficients on greenhouse gases and downstream hydrology and sediment---remain to be resolved \parencite{Mitchell2022}. Carbon budgets should also be integrated into wetland design models when very high nitrate loadings require carbon supplementation, within a more balanced and integrated approach to whole-system wetland evaluation over time \parencite{Kadlec2012}. Dissolved organic matter dynamics likewise demand attention to the range of relevant scales, changes in DOM quality and microbial availability across watershed scales, paired-watershed designs, and sampling that accounts for the episodic influence of subsurface drainage and extreme precipitation \parencite{Dalzell2011,Bellmore2015}.

\subsection{Practice-scale modeling frontiers: wetlands, controlled drainage, and drainage water recycling}

Detailed process modeling of individual conservation practices is a recurring need. Dynamic modeling at time-steps on the order of hours is required to approximate nitrate-N removal in response to the complex flow patterns of tile-drain-fed wetlands, alongside measures to reduce or supplement phosphorus removal and means to buffer and more gradually release tile-drain flows \parencite{Tanner2011}. Constructed-wetland siting and sizing depend on the wetland-to-catchment area ratio and seasonal discharge regime, and long-term studies over a broad range of hydraulic and nutrient loads are needed to establish sustainable nitrogen-retention efficiency \parencite{Steidl2019}. Wetland size ratios need refinement for each local context, together with deeper study of N\textsubscript{2}O emission processes, the N\textsubscript{2}/N\textsubscript{2}O ratio, and pesticide-metabolite accumulation \parencite{Tournebize2017}. Before large created-wetland projects proceed, critical questions of location, design, and function must be explored, supported by adaptive-management research and studies of optimal wetland siting for nitrogen removal across the Mississippi River Basin \parencite{Kovacic2006}. Passing tile-drained water through wetlands is proposed as a leading strategy for lowering nitrogen loads to Midwest rivers and ultimately the Gulf of Mexico, motivating analyses that identify nitrogen-management and drainage-remediation strategies and even reconsideration of river management \parencite{Baeumler2020}. Upstream improvement of water quality before infiltration to aquifers---potentially requiring substantial increases in forested area to lower recharge nitrate---is advanced as a watershed-scale design principle \parencite{Billy2013}. Constructed-wetland studies for irrigation reuse of treated drainage water are ongoing \parencite{Madramootoo2007}, and whole-wetland behavior must be assembled and interpreted over periods long enough to identify temporal trends \parencite{Kadlec2012}.

For controlled and managed drainage, future research should track complete water and nutrient budgets within and across field boundaries \parencite{Shedekar2021} and investigate effects on water and nutrients lost by other pathways---surface runoff, preferential flow, groundwater recharge, and biological uptake \parencite{Ross2016}. Additional research on drainage water management should span a range of soil textures, cropping systems, and climates and evaluate effects on surface and subsurface phosphorus dynamics \parencite{King2022}, with greater attention to phosphorus given the paucity of research and to combinations with practices such as sub-irrigation \parencite{Ross2016}. Whether drainage water management within channelized streams can restore more natural hydrology---and how such hydrologic change affects headwater biogeochemistry---requires comprehensive, integrated hydrology--water-quality--ecology assessment \parencite{King2009}. The potential of controlled tile drainage to reduce downstream flood risk by significantly reducing watershed-scale stream discharge is highlighted for further development \parencite{Sunohara2015}, and control structures on subsurface tile drainage, together with wetland restoration, are noted as promising means of storing more water on the landscape in spring \parencite{Dalzell2018}. Long-term monitoring with moisture sensors and automated systems is sought to improve drainage water recycling efficiency in tile-drained landscapes \parencite{Kaur2021}, and models such as DRAINMOD-DWR are needed to aid the design and management of recycling systems, supported by longer-term performance monitoring across weather, soil, and design conditions and by economic optimization of implementation and management costs \parencite{Moursi2024}. Full-scale field studies of drainage water recycling and saturated buffers under varied conditions, and study of emerging water-quality-wetland types such as reconstructed oxbows, round out the practice-scale agenda \parencite{Mitchell2023}. Restoration of multipurpose oxbows should continue with careful monitoring to determine whether observed relationships with tiles are representative, including further study of temperature differences among oxbow types \parencite{Leberg2024}, and installed weirs and additional floodplain work are being used to resolve surface- and groundwater-flow directions at restored oxbow sites \parencite{Schilling2017}. Edge-of-field filters and phosphorus-sorbing materials warrant testing at higher runoff loadings and in both surface and subsurface pathways \parencite{Plach2023}, filter and grass strips need evaluation of dormant-season loading and vegetation harvest \parencite{Holly2016,Vallée2015}, and the long-term effectiveness and wider deployment of small constructed wetlands and grass strips deserve continued evaluation \parencite{Vallée2015}.

\subsection{Integrated, adaptive, and stakeholder-connected modeling frameworks}

A final cross-cutting frontier is methodological integration. \textcite{Dollinger2015} articulate perhaps the most explicit long-range vision: developing numerically explicit modeling approaches at the landscape scale that integrate and couple the major processes governing the provision of landscape services by ditches, optimizing the spatial arrangement of management practices according to catchment heterogeneity and ditch connectivity, and defining efficient multi-service maximization strategies informed by an a priori, stakeholder-valued ranking of services---while acquiring new knowledge on feedbacks between processes that have most often been studied in isolation. Better mapping of key ditch characteristics (hydraulic, vegetation, and sorption properties) and study of poorly understood processes such as pollutant sorption and degradation are prerequisites \parencite{Dollinger2015,Herzon2008}. Landscape-scale, integrated approaches tailored to local conditions---combining on-farm changes with off-site wetland and riparian restoration and attentive site- and landscape-level design and monitoring---are broadly endorsed \parencite{Blann2009}. Managing catchment nitrogen export is argued to require acknowledging groundwater as well as farm-scale losses and intercepting edge-of-field pathways such as tile and open drains, using stream-rehabilitation and drainage water management tools especially during wet periods \parencite{Goeller2019}. Truly integrated watershed management is said to require better information on the controlling processes---such as spatial groundwater-head distribution---that link topography--geology interactions to the drain fraction \parencite{Mahmood2023}. Coupled field--watershed models enable scenario analyses of water-table management not possible when component models are used separately, and, once tested, can evaluate environmentally sound management in cold regions \parencite{Dayyani2010}. Modeling frameworks should remain extensible---for example, modifying loading-estimation approaches to incorporate additional processes and pathways and the distribution of planting dates across a watershed \parencite{Basu2010}, and refining and dynamically estimating proposed indices through field calibration \parencite{Tao2025}. Detailed, site-specific nutrient-management planning is recommended to handle the interactions between site hydrology and management practices \parencite{Kurz2005}. More catchment-scale research on water quality, wetland carbon sequestration, habitat quality, and cultural and educational services---and on the effects of upgrading existing drainage infrastructure rather than only new installations---is needed to fully populate these integrated frameworks \parencite{Mitchell2022}. Ultimately, tighter interaction among stakeholders, experimentalists, and modelers is identified as essential to make watershed modeling a credible basis for prioritizing drainage water management \parencite{Wang2020b,Kröger2013}.